
This study tested four predictions about scale-dependent error patterns in LLM code judges. Two predictions were supported: (P1) judge-execution agreement increases with scale with diminishing returns (Kruskal-Wallis p=0.021), and (P2) different scales exhibit statistically different FP/FN patterns ($\chi^{2}=45.78$, p=3.$27\times 10^{-8})$. Two predictions were falsified: (P3) ensemble voting degrades rather than improves accuracy ($-8.05$\%, McNemar p=1.$45\times 10^{-6})$, and (P4) unanimous agreement correlates with lower rather than higher accuracy.

The central finding is that scale predicts error type: smaller models over-accept (high FPR), larger models under-accept (high FNR). This asymmetry explains the ensemble failure---the over-accepting majority drowns the accurate minority signal.

Scale selection is error-type selection. Practitioners should choose judges based on which error type is more costly in their application, not by naive ensemble combination.

